\section{Business Recommendations and Impact}
\label{sec:insight}

\subsection{What the models rely on}

Importance is measured by permutation on the held-out window rather than by
impurity on the training set, which inflates high-cardinality features
\citep{strobl2007importance} --- not hypothetical with 74 product categories
(Figure~\ref{fig:importance}). \textbf{Problem 2 has one dominant driver, and it
is actionable.} How late the parcel was carries 0.206 of PR-AUC on its own, how
long it took a further 0.074, the size of the order 0.045; nothing else clears
0.011, and the top two account for \textbf{76.9\%} of total positive importance,
which is why a single depth-limited tree captured 96\% of the achievable signal.
For the customer-experience team that is good news: the driver is something the
business already measures. \textbf{Problem 1A's importances must be read with
care.} No single feature carries much --- the largest,
\code{lead\_payment\_type}, costs 0.034 days of MAE --- because permutation
importance understates correlated predictors such as the promise, the route
histories and the distance; this model has no single point of failure.

\subsection{From a probability to a decision, and when not to act}
\label{sec:cost}

A probability is not a decision. Turning one into the other needs a cost, and
Olist publishes none --- so what follows states its assumptions openly and treats
the conclusions as shapes rather than numbers \citep{elkan2001cost}. The
arithmetic is one line: acting is worth it only where
$\text{precision} > c_{\text{action}} / c_{\text{miss}}$ (Table~\ref{tab:cost}).

\begin{table}[htbp]
\centering
\caption{The decision layer under stated (not measured) costs. Break-even
precision is what decides whether a flag is worth acting on at all.}
\label{tab:cost}
\footnotesize
\begin{tabular}{@{}lrrrl@{}}
\toprule
\textbf{Flag} & \textbf{$c_{\text{action}}$ : $c_{\text{miss}}$} &
\textbf{break-even} & \textbf{best achieved} & \textbf{verdict} \\
\midrule
Low-review risk, at delivery & R\$5 : R\$40  & 0.125 & 0.37 at the optimum & \best{act} \\
Late-order flag, at checkout & R\$12 : R\$40 & 0.300 & 0.28 at best & do not expedite \\
\bottomrule
\end{tabular}
\end{table}

\textbf{The review flag pays, comfortably.} Contacting a customer who was never
going to complain costs perhaps R\$5 of agent time; a 1--2 star review that could
have been headed off costs perhaps R\$40 --- standing in for the damage
\citet{deshpande2023logistics} trace through to seller sales. Break-even
precision is 0.125 and the model reaches 0.369 at its expected-benefit optimum,
flagging 1{,}927 of 12{,}855 reviewed orders (15.0\%) of which 712 do go on to
leave a low review, for a net \textbf{R\$1{,}466 per 1{,}000 delivered orders}
(R\$775 counting only customers who have not yet written, Appendix~\ref{app:timing};
Figure~\ref{fig:threshold}); the margin is wide enough to survive a wrong cost
estimate. Capacity changes how the optimum is used rather than where it is: a
team that can work only 5\% of deliveries still captures R\$1{,}107 of those
R\$1{,}466 (Table~\ref{tab:capacity}), because the top of the list is where the
precision is, so the list is simply worked top-down until the agents run out.

\textbf{The late-order flag does not pay --- for an expensive action.} This is the
more useful finding, and reporting discrimination alone would have missed it.
Expediting freight costs perhaps R\$12 against the same R\$40 of downstream
damage, so break-even precision is 0.300 --- and the derived flag never reaches
it: precision peaks at \textbf{0.28} and falls as volume rises. A model with
ROC-AUC 0.782 and 2.9$\times$ the no-skill floor is genuinely informative and
still not worth paying R\$12 an order to act on. That changes \emph{which} action
it justifies rather than making it useless: break-even falls with the cost of the
intervention, and re-quoting the date and pre-warning the buyer cost almost
nothing.

\subsection{Recommendations by stakeholder}
\label{sec:stakeholders}

\textbf{Operations, at checkout} should work the four levers of
Table~\ref{tab:levers} in order of cost. Every order gets the re-quoted date: the
model's estimate is off by 3.90 days against the promise's 9.79, so the date
shown at checkout can be tightened materially, which \citet{cui2024promise} show
trades directly against returns and retention. Every order above the
recall-oriented threshold of 0.11 gets a proactive notice, which covers five in
six of those that will actually miss. Expedited freight is reserved for a manual
review of the top one or two per cent of the late list, where precision is 0.27,
and that policy is revisited only if expediting falls below about R\$11 an
order. The weekly late series is the monitoring instrument: a week in which
transit lengthens while the promise does not is the February--March 2018
pattern again (Appendix~\ref{app:lateness}).

\textbf{Customer experience, at delivery} should receive each day's deliveries
ranked by $p_{\text{low}}$ and work the list from the top. At ten contacts a day
--- 4.8\% of deliveries --- 57\% of contacts reach a customer who would have left
a 1--2 star review, 35\% one who has not yet written it; twice the capacity
reaches 38\% of them at 44\% precision (Table~\ref{tab:capacity}). The deployed
application returns exactly this list. \textbf{Customer care, on the promised
date,} takes the late parcels: most of their low reviews are written before
delivery, so the customer is reached as the promise breaks, with an apology and
a new date, and the survey is sent after arrival (Appendix~\ref{app:timing}).

\textbf{The platform} should refit both models monthly and monitor the gap
between predicted and realised lead time as the trigger: the drift of
Section~\ref{sec:eda} is the dominant risk to both, and the over-prediction of
Section~\ref{sec:results-segments} is what an unrefitted model looks like after
four months.

\subsection{Measuring what the business case assumes: the outreach pilot}
\label{sec:pilot}

Every reais figure above rests on one assumption the data cannot test: that
contacting a customer who would have left a low review prevents it. Olist never
contacted anyone, so the counterfactual is unobserved, and the evaluators
rightly asked that it be tested rather than assumed. The test is a
randomised controlled experiment \citep{kohavi2009controlled}, and it can be
sized from numbers this report already has (\code{src/pilot\_design.py};
Table~\ref{tab:pilot} and Figure~\ref{fig:pilot}).

\emph{Design.} Each day the model ranks the day's deliveries and the team takes
the top $K$. Within that top $K$, orders are randomised 1:1 --- half contacted,
half left alone --- and the outcome is the 1--2 star rate in each arm.
Randomising \emph{inside} the list is what makes the comparison fair: both arms
have the same predicted risk. The control arm's rate is already known: it is
the model's precision at that daily $K$, 0.572 at ten contacts a day.
\emph{Break-even} is one line of arithmetic: a contact pays when the reduction
in the low-review probability times the review's cost exceeds the contact's
cost, so the required absolute reduction is $5/40 = 0.125$ and the required
relative reduction $0.125/0.572 = 21.9\%$ --- the contact must avert roughly one
low review in five of those it reaches. \emph{Size.} For a two-sided
two-proportion test at $\alpha = 0.05$, detecting a one-in-four reduction
(0.572 to 0.429) with 80\% power needs 191 orders per arm, about \textbf{six
weeks} at ten contacts a day (eight at 90\% power); smaller effects take longer
(Table~\ref{tab:pilot}). The primary outcome is the review rate; the secondary is the
realised benefit per contact, $(\hat p_{\text{control}} - \hat
p_{\text{treated}}) \times \text{R\$}40 - \text{R\$}5$, which replaces the
assumed benefit with a measured one. A second arm splits the late parcels the
same way on their promised date, testing the apology and the later survey.

\subsection{Impact, and what it would take}

Under the stated costs the review-risk list is worth R\$775--1{,}466 per
thousand delivered orders at the cost optimum, depending on whether customers
who have already written count, before the pilot replaces the assumed benefit
with a measured one. The operations levers cost nothing to apply and act on
every order: a promise re-quoted at MAE 3.90 days instead of 9.79, and a
pre-warning that reaches five in six of the orders that will miss. What it
would take is modest: the three artifacts and a daily scoring job, which the
application already runs on a held-out day; a monthly refit; and a six-week
pilot before the customer-experience policy is scaled.
